\documentclass[10pt,a4paper]{amsart}
\usepackage[margin=19mm]{geometry}
\usepackage{amsmath,amssymb,mathtools}
\usepackage[scr=rsfs,cal=euler]{mathalfa}
\usepackage{xfrac}
\usepackage[unicode,hidelinks,bookmarks=false]{hyperref}
\newcommand{\ie}{i.e.\ }
\newcommand{\cf}{cf.\ }
\newcommand{\resp}{respectively\ }
\newcommand{\Subsection}{\subsection}
\providecommand{\nobreakdash}{\nobreak\hbox{-}}
\setlength{\emergencystretch}{2em}
\multlinegap0pt
\usepackage{amssymb}
\usepackage{mathtools}
\usepackage[shortlabels]{enumitem}
\usepackage{xcolor}
\newtheorem{lemma}{Lemma}
\title{Optimality of Wouter van Doorn's Upper Bound for the Mayer--Erdős Farey Problem}
\author{Ricky Cipollini}
\date{Source: arXiv:2607.23302v1 (CC BY 4.0)}
\begin{document}
\maketitle

\noindent\emph{Source and licence.} Optimality of Wouter van Doorn's Upper Bound for the Mayer--Erdős Farey Problem. Version arXiv:2607.23302v1 (CC BY 4.0), distributed by arXiv under the Creative Commons Attribution 4.0 International licence, http://creativecommons.org/licenses/by/4.0/, as recorded in the arXiv licence metadata for this exact version and shown on the abstract page https://arxiv.org/abs/2607.23302v1. Full licence notice: \texttt{licenses/arxiv-2607.23302-CC-BY-4.0.txt}.\par\smallskip

\noindent\textit{Editorial scope.} Counted unit: the article's lemma lem:totient-increments (source0.tex lines 251--264). The article's complete original proof of it is delivered with it (source0.tex lines 266--423, one \texttt{proof} environment of 158 lines). No mathematics is written, repaired or re-derived: the statement and the proof are the article's own lines, in the article's own notation, and no retained line is substituted or re-flowed.  No further source-local context is needed: every cross-reference of every retained line resolves inside the two delivered spans.  Nothing was omitted from any retained span, so the package carries no omission record.  No retained line contains a citation, so the wrapper carries no bibliography and no bibliography entry.  No retained line contains a figure environment, an image-inclusion command or drawing code, so no image is delivered and the package declares no asset dependency.  Editorial formatting only, in addition: an AMS \texttt{amsart} preamble that restates the article's own declarations and macros used by the retained lines, namely the environments \texttt{lemma} on separate counters, together with the standard packages those lines need.  The retained environments print in this document as Lemma 1 in order of appearance, whereas the article prints the same items with the numbering of its own full document.  The complete native source of the article as distributed in the arXiv e-print for this exact version is bundled unchanged as \texttt{sources/206017/source0.tex}.
\begin{lemma}[Totient increments]\label{lem:totient-increments}
For $x>0$ define
\[
  S(x)=\sum_{1\le e<x}\left(1-\frac ex\right)\frac{\varphi(e)}e,
\]
and put $S(0)=0$.  If $x\ge1$ and $y\ge1$, then
\[
  S(x+y)-S(x)\ge \frac y4.
\]
If $x\ge0$ and $y\ge2$, then the same inequality holds.  In particular, for every integer $m\ge2$,
\[
  S(m)\ge \frac m4.
\]
\end{lemma}

\begin{proof}
Let
\[
  \Phi(m)=\sum_{j\le m}\varphi(j).
\]
We first prove the elementary lower bounds
\begin{equation}\label{eq:Phi-first}
  \Phi(m)\ge \frac{m(m+1)}4\qquad(m\ge1)
\end{equation}
and
\begin{equation}\label{eq:Phi-second}
  \Phi(m)\ge \frac{(m+1)^2}{4}\qquad(m\ge7).
\end{equation}
Let
\[
  P(m)=\#\{1\le u,v\le m:(u,v)=1\}.
\]
Counting coprime pairs according to $\max(u,v)$ gives $P(m)=2\Phi(m)-1$.  Indeed, $(1,1)$ contributes $1$, and for each $k\ge2$ the coprime pairs with maximum $k$ contribute $2\varphi(k)$.

If $(u,v)>1$, then some prime $p$ divides both $u$ and $v$.  Hence, by the union bound,
\[
  P(m)\ge m^2-\sum_{p\le m}\left\lfloor\frac mp\right\rfloor^2
        \ge m^2\left(1-\sum_p\frac1{p^2}\right).
\]
We use the explicit estimate
\begin{equation}\label{eq:prime-square-sum}
  \sum_p\frac1{p^2}<\frac{459}{1000}.
\end{equation}
For completeness, here is a verification.  Every prime larger than $97$ is an odd integer at least $99$, and therefore
\[
\sum_p\frac1{p^2}
\le \sum_{p\le97}\frac1{p^2}
   +\sum_{j\ge0}\frac1{(99+2j)^2}.
\]
The tail is bounded by
\[
\sum_{j\ge0}\frac1{(99+2j)^2}
\le \frac1{99^2}+\int_0^\infty\frac{dt}{(99+2t)^2}
=\frac1{99^2}+\frac1{198}<\frac6{1000}.
\]
An exact rational summation over
\[
2,3,5,7,11,13,17,19,23,29,31,37,41,43,47,53,59,61,67,71,73,79,83,89,97
\]
gives $\sum_{p\le97}p^{-2}<453/1000$, proving \eqref{eq:prime-square-sum}.  Thus
\[
  P(m)>\frac{541}{1000}m^2.
\]
For $m\ge24$,
\[
  \frac{541}{1000}m^2\ge\frac{(m+1)^2}{2}-1,
\]
because this is equivalent to $41m^2-1000m+500\ge0$.  Hence, for $m\ge24$,
\[
  2\Phi(m)-1=P(m)>\frac{(m+1)^2}{2}-1,
\]
and so $\Phi(m)>(m+1)^2/4$.

The remaining finite values needed for \eqref{eq:Phi-first} and \eqref{eq:Phi-second} are
\[
\begin{array}{c|rrrrrrrrrrrr}
 m&1&2&3&4&5&6&7&8&9&10&11&12\\ \hline
 \Phi(m)&1&2&4&6&10&12&18&22&28&32&42&46
\end{array}
\]
and
\[
\begin{array}{c|rrrrrrrrrrrr}
 m&13&14&15&16&17&18&19&20&21&22&23&24\\ \hline
 \Phi(m)&58&64&72&80&96&102&120&128&140&150&172&180.
\end{array}
\]
They verify both bounds.

For $x\in[m,m+1]$, $m\ge1$, we have
\[
  S(x)=A_m-\frac{\Phi(m)}x,
  \qquad A_m=\sum_{j\le m}\frac{\varphi(j)}j.
\]
This remains correct at integer endpoints, because the newly appearing term has coefficient zero.  Put
\[
  F(x)=S(x)-\frac x4.
\]
On $[m,m+1]$,
\[
  F(x)=A_m-\frac{\Phi(m)}x-\frac x4,
  \qquad
  F'(x)=\frac{\Phi(m)}{x^2}-\frac14,
  \qquad
  F''(x)=-\frac{2\Phi(m)}{x^3}<0.
\]
So $F$ is concave on each unit interval.  At integers $m\ge1$,
\[
  S(m+1)-S(m)=\frac{\Phi(m)}{m(m+1)},
\]
and therefore, by \eqref{eq:Phi-first},
\[
  F(m+1)-F(m)=\frac{\Phi(m)}{m(m+1)}-\frac14\ge0.
\]
Thus the integer values $F(m)$ are nondecreasing for $m\ge1$.  Furthermore, for $m\ge7$, \eqref{eq:Phi-second} implies throughout $x\in[m,m+1]$ that
\[
  F'(x)\ge \frac{\Phi(m)}{(m+1)^2}-\frac14\ge0.
\]
Hence $F$ is nondecreasing on $[7,\infty)$.

It remains to control small intervals.  Direct substitution gives, for $x\in[m,m+1]$ with $m=1,\ldots,6$,
\[
F(x+1)-F(x)=
\begin{cases}
\dfrac{x^2-3x+4}{4x(x+1)}, & m=1,\\[1.2ex]
\dfrac{5x^2-19x+24}{12x(x+1)}, & m=2,\\[1.2ex]
\dfrac{x^2-7x+16}{4x(x+1)}, & m=3,\\[1.2ex]
\dfrac{11x^2-69x+120}{20x(x+1)}, & m=4,\\[1.2ex]
\dfrac{(x-15)(x-8)}{12x(x+1)}, & m=5,\\[1.2ex]
\dfrac{17x^2-151x+336}{28x(x+1)}, & m=6.
\end{cases}
\]
Each numerator is positive on its indicated interval, so
\begin{equation}\label{eq:F-shift}
  F(x+1)\ge F(x)\qquad(x\ge1).
\end{equation}

We next prove that, for every integer $m\ge0$,
\begin{equation}\label{eq:max-bound}
  \max_{x\in[m,m+1]}F(x)\le F(m+2).
\end{equation}
For $m=0$, we have $S(x)=0$ on $[0,1]$, so $F(x)=-x/4$ and the maximum is $F(0)=0=F(2)$.  For $m=1,3,5$, the maximum occurs at the right endpoint, as follows from the displayed formula for $F'$ and the values $\Phi(1)=1$, $\Phi(3)=4$, and $\Phi(5)=10$.  For $m\ge7$, \eqref{eq:max-bound} follows from monotonicity of $F$ on $[7,\infty)$.  The remaining cases are direct:
\[
  \max_{x\in[2,3]}F(x)=\frac32-\sqrt2<F(4)=\frac16,
\]
\[
  \max_{x\in[4,5]}F(x)=\frac83-\sqrt6<F(6)=\frac3{10},
\]
and
\[
  \max_{x\in[6,7]}F(x)=\frac{19}{5}-2\sqrt3<F(8)=\frac{57}{140}.
\]
This proves \eqref{eq:max-bound}.

Now let $z=x+y$.  First suppose $x\ge0$ and $y\ge2$.  Choose $m\ge0$ with $x\in[m,m+1]$.  Then $z\ge x+2\ge m+2$.  Since the integer values $F(k)$ are nondecreasing for $k\ge1$, and since $m+2\ge2$, every integer endpoint at or after $m+2$ has value at least $F(m+2)$.  On each unit interval $F$ is concave, hence it lies above the chord joining its endpoint values.  Therefore every $t\ge m+2$ satisfies $F(t)\ge F(m+2)$.  Consequently
\[
  F(z)\ge F(m+2)\ge F(x),
\]
where the last inequality is \eqref{eq:max-bound}.  Hence
\[
  S(x+y)-S(x)=\frac y4+F(x+y)-F(x)\ge\frac y4.
\]

It remains to treat $x\ge1$ and $1\le y<2$.  Choose $m\ge1$ with $x\in[m,m+1]$.  If $z\ge m+2$, the previous paragraph gives $F(z)\ge F(x)$.  Otherwise
\[
  z\in[x+1,m+2]\subset[m+1,m+2].
\]
Since $F$ is concave on $[m+1,m+2]$, its minimum on $[x+1,m+2]$ is attained at one of the endpoints.  Thus
\[
  F(z)\ge \min\{F(x+1),F(m+2)\}.
\]
By \eqref{eq:F-shift}, $F(x+1)\ge F(x)$, and by \eqref{eq:max-bound}, $F(m+2)\ge F(x)$.  Hence again $F(z)\ge F(x)$, and the claimed increment inequality follows.  Finally, taking $x=0$ and $y=m\ge2$ gives $S(m)\ge m/4$.
\end{proof}

\end{document}
